\section{Candidate Retrieval}

Following bi-encoder training, the model is used to generate candidate
reference-entity pairs required for cross-encoder training. Since cross-encoders
jointly process a reference and candidate entity, they provide more detailed
interaction modelling but are computationally more expensive than bi-encoders.
Therefore, the cross-encoder is trained on a reduced candidate set rather than
the complete regulatory entity collection.

The trained bi-encoder encodes regulatory references and regulatory entities into
a shared embedding space. For each reference instance, the cosine similarity
between the reference embedding and all candidate entity embeddings is
calculated. The highest-scoring entities are selected as candidate matches for
subsequent cross-encoder training.

The candidate generation process uses the training portion of the augmented
dataset to produce challenging negative examples. For each reference, the known
linked regulatory entity is retained as the positive candidate, while highly
similar but incorrect entities retrieved by the bi-encoder are selected as hard
negative candidates. These candidates represent realistic retrieval errors,
forcing the cross-encoder to learn finer-grained distinctions between similar
regulatory entities.

The generated candidate pairs are transformed into a binary classification
dataset for cross-encoder training. Positive pairs consist of the original
reference paired with its corresponding regulatory entity, while negative pairs
consist of the same reference paired with incorrect candidates retrieved by the
bi-encoder.

The candidate retrieval process therefore bridges the two retrieval stages:
the bi-encoder provides efficient large-scale candidate generation, while the
cross-encoder learns to perform more precise relevance estimation over a smaller
set of potentially relevant regulatory entities.